\documentclass{article}
\usepackage{amsmath, amssymb}
\usepackage[swedish]{babel}

\begin{document}

\section*{Härledning av Vasicek-obligationspriset}

\subsection*{1. Bankkonto och diskonteringsfaktor}

Låt $B_t = $ pengar i konto vid tid $t$. Vid fix ränta $r$ som betalas ut $n$ gånger per år gäller
\[
B_t = \left(1 + \tfrac{r}{n}\right)^{nt} \;\xrightarrow{n \to \infty}\; e^{rt}.
\]
För varierande (stokastisk) ränta $r_s$ generaliseras detta till
\begin{equation}
B_t = e^{\int_0^t r_s\,ds}. \tag{$\star$}
\end{equation}

\subsection*{2. OU-modell för korträntan}

Antag att korträntan följer en Ornstein--Uhlenbeck-process under det riskneutrala måttet $\mathbb{Q}$:
\begin{equation}
dr_t = \theta(\mu - r_t)\,dt + \sigma\,dW_t^{\mathbb{Q}}.
\end{equation}

\subsection*{3. FTAP: obligationspriset som väntevärde}

Enligt finansens fundamentalsats (FTAP) ges priset på en nollkupongsobligation med förfall $T$ av det diskonterade väntevärdet under $\mathbb{Q}$:
\[
P(t,T) = \mathbb{E}^{\mathbb{Q}}\!\left[\, \frac{B_t}{B_T} \;\Big|\; \mathcal{F}_t \,\right]
       = \mathbb{E}^{\mathbb{Q}}\!\left[\, e^{-\int_t^T r_s\,ds} \;\Big|\; \mathcal{F}_t \,\right].
\]
Eftersom $r_t$ är en Markovprocess reduceras betingningen:
\begin{equation}
P(t,T) = \mathbb{E}^{\mathbb{Q}}\!\left[\, e^{-\int_t^T r_s\,ds} \;\Big|\; r_t \,\right].
\end{equation}

\subsection*{4. Feynman--Kac: från väntevärde till PDE}

Definiera
\begin{equation}
u(t,r) = \mathbb{E}^{\mathbb{Q}}\!\left[\, e^{-\int_t^T r_s\,ds} \;\Big|\; r_t = r \,\right],
\end{equation}
så att $P(t,T) = u(t, r_t)$. Feynman--Kac-satsen ger då att $u$ löser
\[
\partial_t u + a(r,t)\,\partial_r u + \tfrac{1}{2} b(r,t)\,\partial_r^2 u - r u = 0,
\qquad u(T,r) = 1,
\]
där koefficienterna avläses från (1):
\[
a(r,t) = \theta(\mu - r), \qquad b(r,t) = \sigma^2.
\]

\subsection*{5. Vasicek-PDE}

Insättning ger obligationsprisets PDE:
\[
\boxed{\;\partial_t P + \theta(\mu - r)\,\partial_r P + \tfrac{1}{2}\sigma^2 \partial_{rr} P - r P = 0, \qquad P(T,T) = 1.\;}
\]

\end{document}